% ============================================================
%  Chapter 14: The Floor Theorem
%  Source: unconstrained-subtask-recursion.tex
% ============================================================

\epigraph{%
  No bounded receiver can know everything.\\
  This is not a limitation.  It is the axiom.%
}{}

\section{The Floor Theorem}
\label{sec:ch14-floor}

\begin{theorem}[Floor Theorem]
  For any receiver $R = (\Sigma_R, \Phi_R, K_R, \mathrm{dec}_R)$ satisfying
  BPSL and any non-trivial question $q$:
  \begin{equation}
    \mathbf{S}_\mathrm{flat}(R) \;\stackrel{\mathrm{def}}{=}\;
    \inf_{K_R \in \Sigma_R} \mathbf{S}[R, q] > 0.
  \end{equation}
  No bounded receiver can attain $S = 0$ on any non-trivial question.
\end{theorem}

\begin{proof}[Proof sketch]
  $\Sigma_R$ is bounded (BPSL clause~i); its partition algebra $\mathcal{A}_\infty$
  is countable (Consequence~\ref{cons:partitionalgebra}).
  The truth about $q$ is a single point in the truth space $\mathcal{T}_q$.
  The receiver's finest resolution is the atoms of $\mathcal{A}_\infty$, which
  are cells of positive measure.
  A positive-measure cell cannot isolate a single point; therefore
  $\mathbf{S}[R,q] \geq \kB\log(\mu(\mathrm{smallest cell})) > 0$.
\end{proof}

\begin{remark}[Physical interpretation]
  The Floor Theorem is the information-theoretic analogue of Heisenberg's
  uncertainty principle: no receiver with finite information capacity can make
  an exact measurement.
  But the Floor Theorem is stronger: it does not depend on quantum mechanics.
  It holds for any bounded receiver, classical or quantum.
  The uncertainty principle is a special case.
\end{remark}

\section{Information Catalysts}
\label{sec:ch14-catalysts}

A surprising consequence of the S-calculus is the existence of
\emph{information catalysts}: observations that reduce the S-value of multiple
questions simultaneously, with multiplicative rather than additive benefit.

\begin{theorem}[Information Catalyst]
  If observation $o$ satisfies
  $\mathbf{S}[R \oplus o, q_1] \cdot \mathbf{S}[R \oplus o, q_2] <
  \mathbf{S}[R,q_1] \cdot \mathbf{S}[R,q_2]$,
  then $o$ is an information catalyst: it reduces the joint S-value
  super-additively.
\end{theorem}

In mass spectrometry, a high-resolution MS/MS spectrum is an information
catalyst: a single fragmentation event simultaneously reduces $\Sk$
(structure), $\St$ (retention time prediction), and $\Se$ (partition depth)
for multiple candidate molecules.

\section{The Unconstrained Subtask Theorem}
\label{sec:ch14-subtask}

\begin{theorem}[Unconstrained Subtask Recursion]
  If the global S-value $\mathbf{S}[R, q]$ is fixed, the local S-values of
  independent subtasks $q_1, q_2, \ldots, q_k$ (where $q = q_1 \wedge \cdots
  \wedge q_k$) are individually unconstrained: any assignment
  $(\mathbf{S}[R,q_1], \ldots, \mathbf{S}[R,q_k])$ consistent with the
  additivity axiom is achievable.
\end{theorem}

This theorem underpins the GPU multi-pass architecture of Chapter~\ref{ch:gpu}:
the global spectral reconstruction task is decomposed into independent shader
passes, each with its own S-value, free to be minimised independently.

%%  SOURCE: integrate full proofs from
%%          unconstrained-subtask-recursion.tex

\bigskip
\begin{tcolorbox}[colback=crownblue!5, colframe=crownblue!60!black,
                  title={\textbf{Chapter 14 Summary}}]
\begin{itemize}[noitemsep]
  \item The Floor Theorem: no bounded receiver attains $S = 0$.
    The floor is strictly positive for any non-trivial question.
  \item Information catalysts reduce multiple S-values simultaneously with
    multiplicative benefit.
  \item The Unconstrained Subtask Theorem: local S-values of independent
    subtasks are free when the global S-value is fixed.
\end{itemize}
\end{tcolorbox}
